% ============================================================
% P2VG — Section III: Methods
% For IEEE Transactions template (IEEEtran class)
% ============================================================

\section{Methods}

\subsection{Overview}

Tổng quan kiến trúc P2VG được trình bày trong Hình~\ref{fig:pipeline}. Cho trước bộ ba thể tích ảnh MRI cột sống --- Sagittal T1-weighted (T1w), Sagittal T2-weighted (T2w) và Axial T2w --- mục tiêu là sinh tự động báo cáo X-quang~$\mathcal{R}$ có độ chính xác lâm sàng cao. Pipeline của chúng tôi gồm bốn giai đoạn: (i)~mô-đun \textbf{3D Wavelet Fusion} kết hợp hai chuỗi sagittal thành một thể tích duy nhất, bảo toàn thông tin chẩn đoán bổ sung từ cả hai contrast; (ii)~hai \textbf{Vision Encoder} song song trích xuất đặc trưng ở mức patch từ ảnh sagittal đã fuse và ảnh axial; (iii)~tầng \textbf{Tri-view Fusion} tích hợp thông tin giữa các mặt phẳng giải phẫu; và (iv)~\textbf{3D Resampler} nén các visual token và chiếu chúng vào không gian nhúng của Mô hình Ngôn ngữ Lớn~(LLM), thực hiện sinh báo cáo thông qua giải mã tự hồi quy.

\begin{figure*}[t]
\centering
\includegraphics[width=0.98\textwidth]{Images/architecture.png}
\caption{Overview of P2VG. Sagittal T1w and T2w are first fused by the 3D Wavelet Fusion module, while axial T2w is processed in parallel. Two ViT3D encoders extract view-specific patch features, which are integrated by the tri-view fusion layer. A 3D resampler compresses features into 256 visual tokens and projects them to the MedGemma embedding space for report generation.}
\label{fig:pipeline}
\end{figure*}

Toàn bộ quá trình sinh báo cáo có thể được mô tả bằng công thức sau:
\begin{equation}
\mathcal{R} = f_{\text{LLM}}\!\Big(\, f_{\Phi}\!\big(\, f_{\text{tri}}\!\big(\, f_{\mathcal{E}}^{\text{sag}}\!\big(\varphi(V^{\text{T1}},\, V^{\text{T2}})\big),\;\, f_{\mathcal{E}}^{\text{ax}}(\tilde{V}^{\text{ax}})\big)\big),\;\, \mathbf{P}\,\Big),
\label{eq:overview}
\end{equation}
trong đó $\varphi$ là bước tiền xử lý wavelet fusion, $f_{\mathcal{E}}^{\text{sag}}$ và $f_{\mathcal{E}}^{\text{ax}}$ là hai bộ mã hóa thị giác kép, $f_{\text{tri}}$ là tầng tri-view fusion, $f_{\Phi}$ là 3D resampler, $f_{\text{LLM}}$ là mô hình ngôn ngữ, và~$\mathbf{P}$ là prompt văn bản.


% ----------------------------------------------------------------
\subsection{3D Wavelet Fusion}

Trong quy trình chụp MRI cột sống lâm sàng, chuỗi Sagittal T1w và T2w cung cấp thông tin chẩn đoán bổ sung cho nhau: ảnh T1w thể hiện rõ hình thái giải phẫu, thành phần mỡ và tín hiệu tủy xương, trong khi ảnh T2w làm nổi bật dịch não tủy, mức độ hydrat hóa đĩa đệm và phù nề --- đây là các đặc điểm quan trọng để nhận diện thoái hóa đĩa đệm, thoát vị và các thay đổi viêm. Thay vì xử lý hai chuỗi như các kênh riêng biệt (sẽ tăng gấp đôi chi phí tính toán của bộ mã hóa thị giác) hoặc đơn thuần lấy trung bình (sẽ mất đi chi tiết đặc trưng theo phương thức), chúng tôi áp dụng chiến lược fuse dựa trên wavelet nhằm lựa chọn và bảo toàn các thành phần thông tin quan trọng nhất từ mỗi chuỗi.

Cho hai thể tích sagittal gốc $V^{\text{T1}},\, V^{\text{T2}} \in \mathbb{R}^{D \times H \times W}$, với độ phân giải lâm sàng thường gặp $15\!\times\!1024\!\times\!1024$, trước tiên chúng tôi chuẩn hóa mỗi thể tích về khoảng $[0, 1]$:
\begin{equation}
\hat{V}(x,y,z) = \frac{V(x,y,z) - V_{\min}}{V_{\max} - V_{\min}},
\label{eq:normalize}
\end{equation}
trong đó $V_{\min} = \min_{x,y,z}\, V(x,y,z)$ và $V_{\max} = \max_{x,y,z}\, V(x,y,z)$.

Quá trình fuse được thực hiện theo từng lát cắt dọc theo trục depth bằng biến đổi wavelet rời rạc 2D~(2D DWT) sử dụng cơ sở Haar~(Daubechies-1). Với mỗi lát cắt tại vị trí depth~$d$, DWT phân rã ảnh thành một băng tần xấp xỉ (approximation, nắm bắt cấu trúc giải phẫu tần số thấp) và ba băng tần chi tiết (detail, nắm bắt thông tin cạnh và kết cấu tần số cao):
\begin{equation}
\big\{\, c_{A,d},\;\; (c_{H,d},\; c_{V,d},\; c_{D,d})\,\big\} = \text{DWT}_{2\text{D}}\!\big(\hat{V}_{:,:,d},\; \text{`haar'}\big),
\label{eq:dwt}
\end{equation}
trong đó $c_A$ là các hệ số xấp xỉ, còn $c_H$, $c_V$, $c_D$ lần lượt là các hệ số chi tiết ngang, dọc và chéo. Phép phân rã này được áp dụng độc lập cho cả $\hat{V}^{\text{T1}}$ và $\hat{V}^{\text{T2}}$, thu được các cặp hệ số băng tần tại mỗi depth.

Quá trình fuse tuân theo hai quy tắc bổ sung, dựa trên nguyên lý rằng cấu trúc giải phẫu được chia sẻ giữa các chuỗi ảnh, trong khi chi tiết bệnh lý có thể xuất hiện rõ ràng hơn ở một chuỗi nhất định:

\textbf{Low-frequency fusion.} Các băng tần xấp xỉ, đại diện cho cấu trúc giải phẫu tổng thể chung của cả hai loại ảnh, được kết hợp bằng phép lấy trung bình nhằm bảo toàn biểu diễn cấu trúc đồng thuận:
\begin{equation}
c_{A,d}^{F} = \frac{c_{A,d}^{\text{T1}} + c_{A,d}^{\text{T2}}}{2}.
\label{eq:lowfreq}
\end{equation}

\textbf{High-frequency fusion.} Các băng tần chi tiết nắm bắt cạnh, ranh giới và mẫu kết cấu thường đặc trưng cho bệnh lý (ví dụ: tín hiệu đĩa đệm sáng trong T2w, viền xương trong T1w). Chúng được fuse bằng quy tắc giá trị tuyệt đối cực đại. Phương pháp lựa chọn dựa trên năng lượng này giữ lại hệ số chi tiết có đáp ứng mạnh hơn tại mỗi vị trí không gian, đảm bảo rằng các đặc trưng có ý nghĩa chẩn đoán từ bất kỳ phương thức nào đều được bảo toàn:
\begin{equation}
c_{H,d}^{F}(i,j) = 
\begin{cases}
c_{H,d}^{\text{T1}}(i,j), & \text{if } \lvert c_{H,d}^{\text{T1}}(i,j) \rvert \geq \lvert c_{H,d}^{\text{T2}}(i,j) \rvert, \\[2pt]
c_{H,d}^{\text{T2}}(i,j), & \text{otherwise}.
\end{cases}
\label{eq:highfreq}
\end{equation}
Quy tắc lựa chọn tương tự được áp dụng cho $c_{V,d}^{F}$ và $c_{D,d}^{F}$. Lát cắt đã fuse sau đó được tái tạo thông qua biến đổi wavelet ngược (inverse DWT):
\begin{equation}
V_{:,:,d}^{F} = \text{IDWT}_{2\text{D}}\!\big(\, c_{A,d}^{F},\;\; (c_{H,d}^{F},\; c_{V,d}^{F},\; c_{D,d}^{F}),\; \text{`haar'}\,\big).
\label{eq:idwt}
\end{equation}

Lặp quá trình trên cho tất cả các lát cắt depth ta thu được thể tích sagittal đã fuse~$V^{F}$. Thể tích này sau đó được tái lấy mẫu (resample) về độ phân giải chuẩn thông qua phép nội suy tuyến tính ba chiều (trilinear interpolation):
\begin{equation}
\tilde{V}^{\text{sag}} = \text{Interpolate}_{\text{trilinear}}\!\big(V^{F},\;\; (D_t, H_t, W_t)\big) \;\in\; \mathbb{R}^{1 \times 32 \times 256 \times 256},
\label{eq:resize}
\end{equation}
trong đó $(H_t, W_t, D_t) = (256, 256, 32)$ là kích thước không gian mục tiêu và chiều dẫn đầu $C\!=\!1$ phản ánh bản chất đơn kênh của thể tích đã fuse. Thể tích Axial T2w cũng trải qua cùng quá trình chuẩn hóa~\eqref{eq:normalize} và tái lấy mẫu một cách độc lập, thu được $\tilde{V}^{\text{ax}} \in \mathbb{R}^{1 \times 32 \times 256 \times 256}$.

Chiến lược wavelet fusion này mang lại hai ưu điểm so với các phương pháp thay thế: (i)~so với đầu vào hai kênh, nó giảm một nửa số lượng token đi vào bộ mã hóa thị giác sagittal, từ đó giảm chi phí tính toán; và (ii)~so với phép lấy trung bình đơn giản, nó bảo toàn thông tin cạnh đặc trưng theo phương thức --- yếu tố quan trọng cho việc nhận diện ranh giới đĩa đệm, gai xương và biên hẹp ống sống.


% ----------------------------------------------------------------
\subsection{Dual Vision Encoder}

Do ảnh MRI theo mặt phẳng sagittal và axial nắm bắt các góc nhìn giải phẫu khác nhau căn bản --- mặt phẳng sagittal thể hiện sự sắp xếp dọc của cột sống, hình thái đĩa đệm tại từng tầng đốt sống, và độ cong lordosis/kyphosis, trong khi mặt phẳng axial cung cấp chi tiết thiết diện ngang của ống sống, lỗ liên hợp thần kinh và khớp facet --- chúng tôi sử dụng hai bộ mã hóa thị giác độc lập, mỗi bộ chuyên biệt trong việc trích xuất đặc trưng theo từng mặt phẳng.

Cả hai bộ mã hóa chia sẻ cùng kiến trúc ViT3D~\cite{dosovitskiy2021image} nhưng duy trì trọng số độc lập, cho phép mỗi bộ học các biểu diễn đặc trưng riêng theo từng view. ViT3D mở rộng kiến trúc Vision Transformer gốc sang đầu vào thể tích 3D thông qua các giai đoạn sau.

\textbf{Patch Embedding.} Thể tích đầu vào $\tilde{V} \in \mathbb{R}^{1 \times 32 \times 256 \times 256}$ được phân chia thành~$N$ patch 3D không chồng lấp có kích thước $(p_D, p_H, p_W) = (4, 16, 16)$. Mỗi patch được duỗi phẳng (flatten) và chiếu tuyến tính vào không gian nhúng $d$ chiều. Số lượng patch là:
\begin{equation}
N = \frac{D_t}{p_D} \times \frac{H_t}{p_H} \times \frac{W_t}{p_W} = \frac{32}{4} \times \frac{256}{16} \times \frac{256}{16} = 2048.
\label{eq:num_patches}
\end{equation}

Một token phân loại (CLS) có thể học $\mathbf{x}_{\text{cls}} \in \mathbb{R}^{d}$ được nối vào đầu chuỗi patch, và các nhúng vị trí có thể học $\mathbf{E}_{\text{pos}} \in \mathbb{R}^{(N+1) \times d}$ được cộng thêm để mã hóa quan hệ không gian:
\begin{equation}
\mathbf{z}_0 = \big[\, \mathbf{x}_{\text{cls}} ;\;\; \mathbf{E}\,\mathbf{p}_1 ;\;\; \mathbf{E}\,\mathbf{p}_2 ;\;\; \cdots ;\;\; \mathbf{E}\,\mathbf{p}_N \,\big] + \mathbf{E}_{\text{pos}},
\label{eq:patch_embed}
\end{equation}
trong đó $\mathbf{p}_i \in \mathbb{R}^{C \cdot p_D \cdot p_H \cdot p_W}$ là patch thứ $i$ đã duỗi phẳng, $\mathbf{E} \in \mathbb{R}^{d \times (C \cdot p_D \cdot p_H \cdot p_W)}$ là ma trận chiếu patch, và $d = 768$ là chiều ẩn (hidden dimension).

\textbf{Transformer Encoding.} Chuỗi nhúng được xử lý qua $L = 12$ khối transformer tiêu chuẩn. Mỗi khối áp dụng multi-head self-attention~(MHSA) theo sau bởi feed-forward network~(FFN), cả hai đều sử dụng pre-normalization và kết nối tắt (residual connection):
\begin{align}
\mathbf{z}'_\ell &= \text{MHSA}\!\big(\text{LN}(\mathbf{z}_{\ell-1})\big) + \mathbf{z}_{\ell-1}, \label{eq:mhsa} \\[4pt]
\mathbf{z}_\ell  &= \text{FFN}\!\big(\text{LN}(\mathbf{z}'_\ell)\big) + \mathbf{z}'_\ell, \label{eq:ffn_block}
\end{align}
trong đó $\text{LN}(\cdot)$ là Layer Normalization. Mô-đun MHSA hoạt động với $h = 12$ attention head và chiều mỗi head $d_k = d/h = 64$. FFN gồm hai phép chiếu tuyến tính với hàm kích hoạt GELU và chiều ẩn $d_{\text{ff}} = 3072$:
\begin{equation}
\text{FFN}(\mathbf{x}) = \mathbf{W}_2\, \text{GELU}\!\big(\mathbf{W}_1 \mathbf{x} + \mathbf{b}_1\big) + \mathbf{b}_2,
\label{eq:ffn}
\end{equation}
trong đó $\mathbf{W}_1 \in \mathbb{R}^{d_{\text{ff}} \times d}$ và $\mathbf{W}_2 \in \mathbb{R}^{d \times d_{\text{ff}}}$.

\textbf{Feature Extraction.} Theo quy ước sử dụng đặc trưng patch cho bài toán hiểu thị giác mật độ cao~\cite{bai2024m3d, chen2024reg2rg}, chúng tôi trích xuất các token patch đầu ra (loại bỏ token CLS) từ tầng cuối cùng đã chuẩn hóa:
\begin{equation}
\mathbf{F} = \text{LN}(\mathbf{z}_L)_{[1{:}N]} \;\in\; \mathbb{R}^{N \times d} = \mathbb{R}^{2048 \times 768}.
\label{eq:feature_extract}
\end{equation}

Áp dụng bộ mã hóa này cho cả hai view, chúng tôi thu được hai tập đặc trưng patch:
\begin{align}
\mathbf{F}_{\text{sag}} &= f_{\mathcal{E}}^{\text{sag}}\!\big(\tilde{V}^{\text{sag}}\big) \;\in\; \mathbb{R}^{B \times 2048 \times 768}, \label{eq:feat_sag} \\[2pt]
\mathbf{F}_{\text{ax}}  &= f_{\mathcal{E}}^{\text{ax}}\!\big(\tilde{V}^{\text{ax}}\big)   \;\in\; \mathbb{R}^{B \times 2048 \times 768}. \label{eq:feat_ax}
\end{align}

Cả hai bộ mã hóa thị giác được khởi tạo từ trọng số ViT3D đã pretrain~\cite{bai2024m3d} và sau đó được fine-tune trong quá trình huấn luyện, cho phép thích ứng miền (domain adaptation) đồng thời tận dụng các biểu diễn đặc trưng 3D tổng quát đã học được trong giai đoạn pretraining.


% ----------------------------------------------------------------
\subsection{Tri-view Fusion}

Bộ mã hóa kép tạo ra hai tập đặc trưng patch đại diện cho hai góc nhìn giải phẫu trực giao. Một cách tiếp cận đơn giản là nối (concatenate) các đặc trưng theo chiều token, nhưng điều này sẽ tăng gấp đôi độ dài chuỗi lên 4096~token, gây ra chi phí tính toán đáng kể cho mô hình ngôn ngữ hạ nguồn. Thay vào đó, chúng tôi fuse đặc trưng theo chiều feature, duy trì số lượng token không gian trong khi cho phép mỗi vị trí token tổng hợp thông tin từ cả hai view.

Cụ thể, đặc trưng patch sagittal và axial được nối theo chiều cuối cùng, sau đó chiếu trở lại chiều đặc trưng ban đầu thông qua một tầng tuyến tính có thể học kết hợp với hàm kích hoạt GELU:
\begin{align}
\mathbf{F}_{\text{cat}} &= \big[\, \mathbf{F}_{\text{sag}} ;\;\; \mathbf{F}_{\text{ax}} \,\big] \;\in\; \mathbb{R}^{B \times 2048 \times 1536}, \label{eq:concat} \\[4pt]
\mathbf{F}_{\text{tri}} &= \text{GELU}\!\big(\mathbf{W}_f \, \mathbf{F}_{\text{cat}} + \mathbf{b}_f\big) \;\in\; \mathbb{R}^{B \times 2048 \times 768}, \label{eq:trifuse}
\end{align}
trong đó $\mathbf{W}_f \in \mathbb{R}^{768 \times 1536}$ và $\mathbf{b}_f \in \mathbb{R}^{768}$ là các tham số có thể học. Thiết kế này cho phép mô hình học các tương tác bổ sung giữa hai view tại mỗi vị trí không gian --- ví dụ, liên kết hình thái đĩa đệm quan sát được trên mặt phẳng sagittal với đường kính ống sống tương ứng đo được trên mặt phẳng axial --- đồng thời giữ chuỗi token ở độ dài 2048 có thể quản lý được.


% ----------------------------------------------------------------
\subsection{3D Resampler}

Mặc dù tri-view fusion tạo ra các biểu diễn thị giác phong phú, số lượng 2048~token sẽ tạo gánh nặng lớn cho cơ chế self-attention của mô hình ngôn ngữ, vốn có độ phức tạp tăng theo bình phương với độ dài chuỗi. Để giải quyết vấn đề này, chúng tôi đề xuất 3D Resampler nhằm nén các visual token thông qua spatial pooling và chiếu chúng vào không gian nhúng của LLM. Thành phần này gồm hai giai đoạn: spatial pooling và dimensional projection.

\textbf{Spatial Pooling.} Chuỗi đặc trưng đã fuse được reshape về lưới không gian 3D ban đầu và nén thông qua 3D average pooling:
\begin{align}
\mathbf{F}_{\text{3d}} &= \text{Reshape}\!\big(\mathbf{F}_{\text{tri}},\;\; (B,\, 768,\, 8,\, 16,\, 16)\big), \label{eq:reshape3d} \\[4pt]
\mathbf{F}_{\text{pool}} &= \text{AvgPool3D}\!\big(\mathbf{F}_{\text{3d}},\;\; k\!=\!2,\; s\!=\!2\big) \;\in\; \mathbb{R}^{B \times 768 \times 4 \times 8 \times 8}, \label{eq:avgpool} \\[4pt]
\mathbf{F}_{\text{seq}} &= \text{Reshape}\!\big(\mathbf{F}_{\text{pool}},\;\; (B,\, N_v,\, 768)\big), \label{eq:reshape_seq}
\end{align}
trong đó $N_v = 4 \times 8 \times 8 = 256$. Phép biến đổi này giảm số lượng visual token từ 2048 xuống 256 --- nén $8\times$ giúp giảm đáng kể chi phí tính toán của mô hình ngôn ngữ trong khi vẫn bảo toàn tính cục bộ không gian. Việc sử dụng pooling 3D (thay vì 1D) đảm bảo rằng quá trình nén tôn trọng cấu trúc không gian thể tích gốc, tránh mất tương ứng giải phẫu vốn xảy ra khi giảm mẫu token một cách thô sơ.

\textbf{MedGemma-Pretrained Projection.} Để thu hẹp khoảng cách chiều giữa bộ mã hóa thị giác ($d = 768$) và mô hình ngôn ngữ ($d_{\text{LLM}} = 2560$), chúng tôi áp dụng thiết kế projection lấy cảm hứng từ kiến trúc đa phương thức MedGemma~\cite{medgemma}. Mô-đun này gồm một adapter có thể huấn luyện, một RMSNorm đã pretrain, và một phép chiếu tuyến tính đã pretrain:
\begin{align}
\mathbf{h}  &= \text{GELU}\!\big(\mathbf{W}_a \, \mathbf{F}_{\text{seq}} + \mathbf{b}_a\big) \;\in\; \mathbb{R}^{B \times 256 \times 1152}, \label{eq:adapter} \\[4pt]
\hat{\mathbf{h}} &= \text{RMSNorm}\!\big(\mathbf{h}\big) \;\in\; \mathbb{R}^{B \times 256 \times 1152}, \label{eq:rmsnorm} \\[4pt]
\mathbf{F}_{\text{vis}} &= \mathbf{W}_p \, \hat{\mathbf{h}} \;\in\; \mathbb{R}^{B \times 256 \times 2560}, \label{eq:projection}
\end{align}
trong đó $\mathbf{W}_a \in \mathbb{R}^{1152 \times 768}$ là adapter có thể huấn luyện, ánh xạ đặc trưng ViT3D tùy chỉnh của chúng tôi vào không gian đầu vào 1152 chiều mà MedGemma kỳ vọng. Tầng RMSNorm (với tham số scale $\boldsymbol{\gamma} \in \mathbb{R}^{1152}$) và phép chiếu $\mathbf{W}_p \in \mathbb{R}^{2560 \times 1152}$ tái sử dụng trọng số đã pretrain từ visual projector gốc của MedGemma và được giữ cố định trong quá trình huấn luyện. Các thành phần pretrain này mã hóa sự căn chỉnh thị giác-ngôn ngữ bền vững, đã được học từ các cặp ảnh y tế-văn bản quy mô lớn. Thiết kế này cho phép framework của chúng tôi tận dụng pretraining y tế của MedGemma trong khi chỉ cần huấn luyện adapter nhẹ từ đầu.


% ----------------------------------------------------------------
\subsection{Large Language Model}

Giai đoạn cuối cùng sử dụng một mô hình ngôn ngữ lớn đã pretrain để sinh báo cáo X-quang, có điều kiện trên các đặc trưng thị giác và prompt văn bản. Chúng tôi sử dụng backbone văn bản của MedGemma 1.5 4B-IT, dựa trên kiến trúc Gemma~3 --- một transformer chỉ có decoder với 26~tầng và chiều ẩn~2560.

\textbf{Multimodal Prompt Construction.} Theo định dạng hướng dẫn của Gemma~3, các visual token được xen kẽ với prompt văn bản để tạo thành đầu vào đa phương thức. Chúng tôi giới thiệu token đặc biệt \texttt{<im\_patch>} để đánh dấu vị trí visual token. Trong quá trình forward pass, 256~nhúng placeholder được thay thế bằng các đặc trưng thị giác đã chiếu~$\mathbf{F}_{\text{vis}}$, tạo thành chuỗi nhúng đa phương thức thống nhất:
\begin{equation}
\mathbf{E}_{\text{mm}} = \big[\, \mathbf{e}_{\text{bos}} ;\;\; \mathbf{F}_{\text{vis}} ;\;\; \mathbf{e}_{q_1} ;\;\; \cdots ;\;\; \mathbf{e}_{q_m} \,\big] \;\in\; \mathbb{R}^{B \times M \times 2560},
\label{eq:multimodal_embed}
\end{equation}
trong đó $\mathbf{e}_{\text{bos}}$ là nhúng đầu chuỗi (beginning-of-sequence) và $\mathbf{e}_{q_i}$ là các nhúng token câu hỏi/prompt.

\textbf{LoRA Fine-Tuning.} Việc fine-tune trực tiếp toàn bộ 4~tỷ tham số của mô hình ngôn ngữ sẽ tốn kém về mặt tính toán và dễ bị overfitting trên tập dữ liệu y tế hạn chế của chúng tôi. Thay vào đó, chúng tôi áp dụng Low-Rank Adaptation~(LoRA)~\cite{hu2022lora}, chèn các ma trận hạng thấp có thể huấn luyện vào mỗi ma trận trọng số cố định trong backbone LLM:
\begin{equation}
\mathbf{W} = \mathbf{W}_0 + \frac{\alpha}{r}\, \mathbf{B}\, \mathbf{A},
\label{eq:lora}
\end{equation}
trong đó $\mathbf{W}_0 \in \mathbb{R}^{d_{\text{out}} \times d_{\text{in}}}$ là trọng số pretrain cố định, $\mathbf{A} \in \mathbb{R}^{r \times d_{\text{in}}}$ và $\mathbf{B} \in \mathbb{R}^{d_{\text{out}} \times r}$ là các ma trận phân rã hạng thấp có thể huấn luyện, $r = 64$ là hạng, và $\alpha = 128$ là hệ số tỷ lệ. LoRA được áp dụng cho tất cả các tầng tuyến tính trong backbone transformer, trong khi vision tower, projector và các tầng embedding được loại trừ khỏi các module mục tiêu LoRA và được huấn luyện toàn bộ gradient.

Chiến lược huấn luyện có chọn lọc này dẫn đến cấu hình ở mức thành phần như được tóm tắt trong Bảng~\ref{tab:trainable}:

\begin{table}[t]
\centering
\caption{Trạng thái huấn luyện của từng thành phần trong P2VG.}
\label{tab:trainable}
\begin{tabular}{ll}
\toprule
\textbf{Thành phần} & \textbf{Trạng thái huấn luyện} \\
\midrule
ViT3D Encoders (sag + ax) & Fine-tune toàn bộ \\
Tri-view Fusion layer      & Fine-tune toàn bộ \\
Projector Adapter ($\mathbf{W}_a$) & Fine-tune toàn bộ \\
MedGemma Projection ($\mathbf{W}_p$, $\boldsymbol{\gamma}$) & Pretrained, cố định \\
LLM Backbone               & Cố định + LoRA adapters \\
New Token Embeddings        & Fine-tune toàn bộ \\
\bottomrule
\end{tabular}
\end{table}


% ----------------------------------------------------------------
\subsection{Training Objective}

Mô hình được huấn luyện đầu-cuối (end-to-end) bằng cách tối thiểu hóa hàm mất mát cross-entropy tự hồi quy trên các token báo cáo. Cho chuỗi nhúng đa phương thức đầu vào~$\mathbf{E}_{\text{mm}}$ và báo cáo ground-truth $\mathcal{R} = \{r_1, r_2, \ldots, r_T\}$, hàm mất mát được tính như sau:
\begin{equation}
\mathcal{L} = -\sum_{t=1}^{T} \log\, p_\theta\!\big(r_t \mid \mathbf{E}_{\text{mm}},\; r_1,\; \ldots,\; r_{t-1}\big),
\label{eq:loss}
\end{equation}
trong đó $\theta$ bao gồm tất cả các tham số có thể huấn luyện. Để tránh việc mô hình được thưởng khi ``dự đoán'' các visual token và question token (vốn được cung cấp sẵn làm đầu vào), hàm mất mát được che (mask) sao cho chỉ có các token câu trả lời đóng góp vào gradient:
\begin{equation}
\ell_i = 
\begin{cases}
-100,       & \text{if } i \leq Q \quad \text{(prompt tokens)}, \\[2pt]
r_{i-Q},    & \text{if } Q < i \leq Q\!+\!T \quad \text{(answer tokens)}, \\[2pt]
-100,       & \text{if } i > Q\!+\!T \quad \text{(padding tokens)},
\end{cases}
\label{eq:label_mask}
\end{equation}
trong đó $Q$ là tổng độ dài phần prompt (bao gồm cả visual token). Việc gán nhãn bằng $-100$ tuân theo quy ước chuẩn của PyTorch để bỏ qua token trong phép tính cross-entropy.


% ----------------------------------------------------------------
\subsection{Inference}

Ở giai đoạn suy luận, mô hình sinh báo cáo theo phương pháp tự hồi quy sử dụng nucleus sampling với top-$p = 0.9$ và nhiệt độ $\tau = 0.7$. Quá trình sinh dừng khi mô hình tạo ra token kết thúc chuỗi (end-of-sequence) hoặc đạt tối đa 256~token mới. Prompt câu hỏi được lấy mẫu ngẫu nhiên từ một tập hợp các mẫu truy vấn lâm sàng được định nghĩa trước (ví dụ: ``Mô tả các bất thường trong ảnh MRI cột sống này'', ``Các phát hiện chính trong ảnh chụp này là gì?'') nhằm khuyến khích tính đa dạng và tính bền vững của kết quả sinh.
